\section{Appendix: Scope and Qualifications}

This study concerns the universal-calendar celebration of the Resumed Sixth Sunday after the Epiphany on 15 November 2026 in the 1962 Roman Missal. The second-class Sunday has its appointed Trinity Preface, Gloria and Creed. Under the governing occurrence rules it has no ordinary commemoration of Saint Albert the Great. No local patronal, institutional or diocesan addition is included. The dated institutional ordos corroborate the Sunday identity; their particular additions do not enlarge this universal inventory.\footnote{\textit{Rubricae generales} 18(b), 91, 109, 111(b); ICRSP France, November 2026 Ordo, entry for 15 November, \url{https://icrspfrance.fr/ordo.php?d=2026-11-01}; FSSPX Ordo, 15 November 2026, \url{https://ordo.fsspx.org/day/2026-11-15?lang=en}.}

All eleven appointed Latin texts are included, retaining the Missal's abbreviated Ordinary cues and short conclusions. The preexisting wording was collated against the 1962 controller and the identified 1862 antecedent. Accents, ligatures and line divisions are normalized. The 1962 book's erroneous Epistle citation is corrected by its actual body and the corresponding winter lesson. The English biblical counterparts are the registered Douay--Rheims (Challoner) text. They are visibly distinct from the Introit's composite adaptation, the shortened Psalm petitions, liturgical introductions and repetitions. The prayers and assembled Preface quote Cummiskey's 1861 historical English, including its own abbreviations; they are not an approved English translation of the 1962 Missal. The historical Secret has no closing cue in that witness, and none is added here.

The historical Latin antecedents, historical English translations and historical patristic translations used here are public-domain material in the United States under their recorded source bases. This does not place the entire modern facsimile deliveries, website revisions or editorial packaging in the public domain. Aquinas is explained in original commentary from the identified Latin delivery, without reproducing its modern editorial text. The accompanying text-collation and source records identify the exact witnesses, rights boundaries and limits of the research.

The two whole-Mass interpretations are source-grounded editorial syntheses. The Fathers are credited with their explanations of particular passages. Schuster directly comments on the three prayers and the two lessons, while his treatment of the matching chants comes from a Mass with other lessons and prayers. His wider historical and political generalizations are not premises of this study. The bounded research established no direct patristic exposition of the three exact orations or this exact Preface. Their own words and the stated later liturgical commentary support their treatment. Augustine's exposition of the full \textit{De profundis} extends beyond the verses appointed in the chants; Bede's treatment of the Communion concerns the surrounding Mark 11:22--25 sequence, without a separately displayed comment headed 11:24.

The chronology appendix distinguishes claims about events, traditional attributions, composition and proposed alternatives. Those dated authorities are presented with their limits; their juxtaposition is not a claim to contemporary scholarly unanimity. This is a study aid, not an official liturgical book or an ecclesiastically approved commentary.

\section{References}

\begin{itemize}
\item \textbf{The appointed texts and their historical witnesses.}
\textit{Missale Romanum}, Vatican typical edition, 1962, pp.~421--422, nos.~1759--1768; p.~293, no.~1082; p.~49, no.~338. Governing rubrics: \textit{Rubricae generales} 18(b), 91, 109, 111(b), and \textit{Rubricae generales Missalis romani} 298. Controller: \url{https://media.churchmusicassociation.org/pdf/missale62.pdf}.
\textit{Missale Romanum}, Ratisbon: Pustet, 1862, pp.~44--45, 239--240, 357--358, supplies the antecedent Latin checked for the appointed bodies. Eugene Cummiskey, publisher, \textit{The Roman Missal Translated into the English Language for the Use of the Laity}, Philadelphia, 1861: Sixth Sunday after Epiphany, printed p.~102 and the subsequent page anomalously numbered 101; Preface material, pp.~xxvii--xxix. Bible: Douay--Rheims (Challoner), Project Gutenberg no.~1581, the registered biblical counterparts reproduced above and their cited contexts. Psalm numbers in the commentary are Vulgate, followed by Hebrew numbers where useful.

\item \textbf{John Chrysostom.}
\textit{Homilies on Matthew}, 46.2--3 and 47.1, historical English in \textit{Nicene and Post-Nicene Fathers}, first series, vol.~10, ed. Philip Schaff, CCEL complete-text witness. \textit{Homilies on First Thessalonians}, homily 1 on 1:2--7, historical English in NPNF first series, vol.~13; the volume's preface attributes the Thessalonian translation to James Tweed. The exposition uses the identified historical English, without claiming a fresh Greek collation.

\item \textbf{Jerome.}
\textit{Commentary on Matthew}, book II, on 13:31--35, in J.-P. Migne, \textit{Patrologia Latina} 26 (Paris, 1845), cols.~89--93; Latin page-image witness. \textit{Commentary on Jeremiah}, book V, on 29:10--14, PL 24 (Paris, 1865), cols.~893--894. The descriptions of his arguments are original English exposition of the checked Latin.

\item \textbf{Augustine.}
\textit{Expositions of the Psalms} (\textit{Enarrationes in Psalmos}), 43.6--8, 84.2--6 and 129.1--8, historical English in NPNF first series, vol.~8, ed. Philip Schaff, CCEL complete-text witness. Its headings are Psalms 44, 85 and 130 respectively; the Latin numbering is used in the commentary's citations. The discussion of the morning watch follows Augustine's lemma and exposition.

\item \textbf{Thomas Aquinas.}
\textit{Super I Epistolam ad Thessalonicenses lectura}, chapter 1, lecture 1 [87862], Latin text of the Turin 1953 edition in the Corpus Thomisticum delivery revised by Enrique Alarcón, \url{https://www.corpusthomisticum.org/c1t.html}. Used by attributed paraphrase for the thanksgiving, active virtues, explanations of power, imitation and expectation of the Son.

\item \textbf{Bede.}
\textit{In Marci evangelium expositio}, on Mark 11:22--25, the Latin Migne-derived Wikisource witness, paragraphs beginning \textit{Et respondens Jesus} and \textit{Et cum stabitis ad orandum}. The retained application concerns faith, charity, forgiveness and eternal life in this connected exposition.

\item \textbf{Robert Bellarmine.}
\textit{Commentary on the Book of Psalms}, trans. John O'Sullivan, Dublin: James Duffy, 1866, historical-text transcription delivered by eCatholic2000: Psalm 43, explanations 7--9; Psalm 84, explanation 1; Psalm 129, explanations 3--5. The explanation numbers are the commentary's, not silently renumbered biblical verses. O'Sullivan's translation omits some philological material from the Latin work.

\item \textbf{Ildefonso Schuster.}
\textit{The Sacramentary}, trans. Arthur Levelis-Marke, London: Burns Oates \& Washbourne, vol.~I (1924), ``Sixth Sunday after the Epiphany,'' pp.~416--417, with the Collect at 416 and the Secret and Postcommunion at 417; vol.~III (1927), pp.~187--188, on the matching Gradual, Alleluia and Offertory. Exact facsimile and passage identities are retained in the source records.
\end{itemize}
